\subsection{Constructing the local representation and period families}
\label{archimedean-pieri-family}

The representation hypotheses in the preceding local norm estimates
can be supplied by classical representation theory. We use
Harish--Chandra's discrete-series classification, the tensor identity
and good-range theorem for cohomological induction, and the
Casselman--Wallach equivalence
\cite{pieri-family-hks,pieri-family-kv,pieri-family-cw}.
These are cited theorems, not formal-source assumptions.
All maps below are continuous on smooth representation spaces;
no boundedness of the tensor projections in the ambient Hilbert norm
is asserted. Throughout the original local consumer has $k\ge1$,
as required by both its numerical family predicate and
$\mathrm{GapOK}$; the case $k=0$ is outside this obligation.

\begin{theorem}\label{original-signed-pieri-family}
For the original numerical family, choose a discrete series
$H_\nu$ of parameter $\lambda_0+\nu$ at every
$\|\nu\|_1\le q_{\rm fam}$. The choice may equal the prescribed
representation at its given index.
The exterior and dual-standard decompositions used by the local
norm argument exist on their smooth spaces, with their exact
parameter shifts and with signed isometric inclusions.
Their inclusions and projections preserve the span of group
translates of compact-finite vectors.
\end{theorem}
\begin{proof}
The numerical family gives the correct integral lattice and gaps
at least two at every permitted index. Harish--Chandra's
classification supplies these representations and identifies
the prescribed one by its parameter and signature.

Fix a parameter $\mu$. Use the $\theta$-stable Borel determined by
its ordered signature. In infinitesimal-character normalization,
its cohomologically induced module is the discrete series at $\mu$.
The finite-dimensional module $\bigwedge^\ell\mathbb C^n$ has
one-dimensional weights $\mathbf1_I$, $|I|=\ell$;
the dual standard module has weights $-e_s$.
Filter either module by this Borel. The tensor identity gives a
filtration with the induced modules at $\mu+\mathbf1_I$,
or $\mu-e_s$, as its factors. All shifted parameters remain in
the same open chamber. The good-range vanishing theorem therefore
concentrates the induction in its one relevant degree.
The shifted parameters are distinct and ordered, so their
infinitesimal characters are distinct. The finite filtration splits
by generalized central-character decomposition. Every factor is
the indicated irreducible discrete series.

Casselman--Wallach globalizes these isomorphisms to smooth
representations, giving precisely the continuous smooth maps required
by the original tensor interface. Pull back the tensor product of
the finite-dimensional invariant Hermitian form and the unitary
Hilbert form. Different summands are orthogonal: the resulting
invariant sesquilinear pairing would otherwise intertwine distinct
irreducible Harish--Chandra modules with their Hermitian duals.
On one summand Schur's lemma makes the form a real scalar multiple
of its unitary form. That scalar is nonzero. Otherwise that summand,
being orthogonal to all the others, would lie in the radical of the
nondegenerate tensor form. Rescale its inclusion by the reciprocal
square root of the scalar's absolute value and rescale the
projection inversely. The scalar becomes $+1$ or $-1$.
The identities extend from compact-finite vectors to smooth vectors
by continuity and density.

An equivariant map sends compact-finite tensors to compact-finite
tensors. For a group translate, intertwining expresses each tensor
coordinate as a finite sum of translates of such vectors.
This proves the last assertion.
\end{proof}

The following argument supplies the smooth coefficient estimate needed
to construct the period maps. It uses auxiliary local representations
outside the numerical family only to prove convergence; none is
substituted for an original arithmetic representation.

\begin{lemma}\label{original-stable-tensor-power}
Let $\mu$ be a regular integral parameter whose noncompact root
gaps are at least $h>0$. For every integer $M$ with $0\le M<h$,
$(\mathbb C^n)^{\otimes M}\otimes H_\mu^\infty$ is a finite
direct sum of smooth discrete series, with a signed orthogonal
decomposition for its original tensor Hermitian form.
Repeated final parameters are allowed.
\end{lemma}
\begin{proof}
Iterate the one-step standard tensor identity, separately on every
summand already constructed. A regular integral parameter has
adjacent gaps at least one. Adding $e_s$ can therefore leave the
open chamber only by an adjacent equality. After at most $M$
steps, every noncompact gap is at least $h-M>0$, so an equality
can involve only a compact simple root.

Such a compact-wall factor is zero. Indeed, induction in stages
through its compact rank-one Levi reduces it to the
Borel--Weil--Bott calculation on $\mathbb P^1$ at degree $-1$:
the infinitesimal-character pairing is zero, so the inducing
weight pairs to $-1$ after subtracting $\rho$.
Both cohomology groups of $\mathcal O_{\mathbb P^1}(-1)$ vanish.
All other factors are regular discrete series in the same chamber.
The one-step central characters are distinct, so that step splits,
as in the proof of Proposition~5.5 of
\cite{pieri-family-hks}. Normalize its Hermitian forms as above.

Tensoring distributes over the previously obtained direct sum.
Thus the induction constructs an actual direct sum labelled by
paths. Equal final parameters in different paths do not require
any assertion that self-extensions split: the relevant summands
were already split at the preceding step. Their orthogonality is
inherited from that step. This also proves the Hermitian assertion.
\end{proof}

\begin{theorem}\label{original-smooth-decay-from-tensor}
Let $N(g)=(\sum_{i,j}|g_{ij}|^2)^{1/2}$ be the exact source norm.
Suppose the actual noncompact gap is at least $h_*$ and the
original condition $\mathrm{GapOK}(k,r,h_*)$ holds. Put
$a=h_* /(8\sqrt{2r+2})$.
There are a constant $C$ and a finite sum $p$ of the original
derived Hilbert seminorms such that
\[
 |\langle\sigma(g)v,w\rangle|
 \le C\,p(v)p(w)N(g)^{1-a}
 \quad(v,w\in H_\mu^\infty,\ g\in G).
\]
No uniform bound on $C$ or on the order of $p$ is needed here.
\end{theorem}
\begin{proof}
Set $M=\lceil a-1\rceil$. The original gap gives $a>1$,
and $M\le a\le h_*/8<h_*$. Apply
Lemma~\ref{original-stable-tensor-power}.
Each projection, being continuous from the smooth space to the
Hilbert space of its summand, is bounded by a finite sum of
derived Hilbert seminorms. There are finitely many projections
and finite-factor basis vectors, so one common such sum $p$
and a finite constant suffice.

Use the signed tensor decomposition to expand a matrix coefficient
between two finite-factor basis vectors tensored with $v,w$.
Unitarity on each summand bounds its absolute value by
$Cp(v)p(w)$. The original finite-factor Hermitian matrix is a
signed permutation matrix. Its tensor powers are likewise.
Consequently this estimate applies to
$g_{ij}^M\langle\sigma(g)v,w\rangle$ for every $i,j$.
Since $N(g)\le n\max_{i,j}|g_{ij}|$,
\[
 N(g)^M|\langle\sigma(g)v,w\rangle|\le n^MCp(v)p(w).
\]
Absorb $n^M$ in the constant. On $G$, $N(g)\ge1$, and
$M\ge a-1$, proving the assertion.
In particular, fixing a compact-finite $w$ gives the exact
existential smooth-seminorm estimate denoted $\mathrm{DecInf}$
in the source. A uniform bound for these auxiliary seminorms
has not been inserted into any later quantitative estimate.
\end{proof}

\begin{theorem}\label{original-normalized-bessel-family}
For every representation in the original numerical family and each
original Bessel side, there exists a continuous smooth Bessel map,
possibly zero, normalized by the exact original subgroup period.
The signed Pieri scalar and adjoint identities then hold.
\end{theorem}
\begin{proof}
Write $S=U(W_0)N$ and
$\lambda(hn)=\psi(n)\pi_W(h)$.
The original compact and nilpotent measures give its Haar measure
$ds$, unchanged. It is a closed unimodular amenable subgroup.
The ambient-volume packing proof of
Section~\ref{archimedean-quotient-mass} applies verbatim to $S$:
left multiplication has ambient real determinant one, so
$\operatorname{vol}\{s:N(s)\le R\}\le C_SR^{2n^2}$.
For the original numerical family the common choice $h_*=h/4$
is permitted by its shift reserve and its stated
$\mathrm{GapOK}(k,r,h/4)$ condition. The last inequality in
$\mathrm{GapOK}$ gives
$a-1\ge9(r+1)n^2+3\,\mathrm{AG}(k,r)+2>2n^2$.
The preceding theorem therefore makes every period absolutely
convergent and continuous in both smooth inputs.

Here use Hermitian forms linear in the first variable. On
$E=H_\mu^\infty\otimes W^*$ consider the unitary $S$-representation
$\tau=\sigma\otimes\lambda^\vee$ and its period form
\[
 Q(x,y)=\int_S\langle\tau(s)x,y\rangle\,ds .
\]
It is positive semidefinite. To see the sign with the original
Haar mass, take a F{\o}lner net $A$ of finite positive measure.
Then
\[
 \int_S\langle\tau(s)x,x\rangle
       \frac{\operatorname{vol}(A\cap sA)}{\operatorname{vol}(A)}
       \,ds
 =\frac1{\operatorname{vol}(A)}
       \left\|\int_A\tau(s^{-1})x\,ds\right\|^2\ge0 .
\]
The overlap ratio tends to one and is bounded by one.
Absolute convergence and dominated convergence give $Q(x,x)\ge0$.

Each $Q(\,\cdot\,,y)$ is a continuous invariant functional.
The original group is $U(k+r,r+1)$, its compact smaller group is
$U(k)$, and its nonzero flag character is generic in exactly the
sense of Jiang--Sun--Zhu. Their Theorem A
\cite{pieri-family-jsz} gives dimension at most one for these
functionals. Thus $Q$ has rank at most one.
If $Q=0$, take $B=0$. Otherwise choose $x_0$ with
$Q(x_0,x_0)>0$ and put
$f(x)=Q(x,x_0)/\sqrt{Q(x_0,x_0)}$.
Positivity and rank one give
$Q(x,y)=f(x)\overline{f(y)}$.
Since $W$ is finite dimensional, there is a unique continuous
map $B:H_\mu^\infty\to W$ with
$f(v\otimes\xi)=\xi(Bv)$.
Invariance gives $B\sigma(s)=\lambda(s)B$.
Reading the matrix entries of $Q$ proves the original operator
period identity
\[
 \int_S\langle\sigma(s)v,w\rangle\lambda(s)^{-1}\,ds
   =R(Bv,Bw),\qquad R(u,z)\eta=\langle\eta,z\rangle u .
\]
In particular no Haar factor, dimension factor or positive
normalizing constant was discarded.

Write $A_Iv=p_I(\omega\otimes v)$ and
$C_Iy=\operatorname{ev}_{\Gamma}(\omega,i_Iy)$.
The original flag vector is $S$-fixed: for an exterior power this
is the wedge of an initial flag segment, and for the dual standard
module it is the flat of the fixed first flag vector.
The signed tensor pairing and the unchanged subgroup period give,
before any scalar is chosen,
\[
 R(BC_Iy,Bv)=\epsilon_I R(B_Iy,B_IA_Iv).
\]
This is the period comparison in
Lemma~\ref{lowering-normalized-adjoint}; absolute convergence
here makes it valid for all smooth $v,y$.
If $B=0$ and $B_I\ne0$, choose $y$ with $B_Iy\ne0$.
The displayed rank-one identity forces $B_IA_Iv=0$ for every $v$.
If both maps are zero the same conclusion is immediate.
Thus $\kappa_I=0$ supplies both the scalar and signed adjoint
identities whenever $B=0$. This argument does not infer the
absence of other Bessel functionals from a zero period form.

If $B\ne0$, continuity and density make it nonzero on the original
period test domain. Multiplicity one gives
$B_IA_Iv=\kappa_I Bv$. Substitute this equality into the displayed
period comparison and choose $v$ with $Bv\ne0$. It gives
$BC_I=\epsilon_I\overline{\kappa_I}B_I$, including $B_I=0$,
when $\kappa_I=0$.
The tensor maps preserve that lemma's test domain by
Theorem~\ref{original-signed-pieri-family}.
The full Capelli character is the given discrete-series
infinitesimal character; differentiating $B\sigma(h)=\pi_W(h)B$
gives its compact Capelli character. Hence the exact
normalized-lowering theorem applies with its stated shifts.
\end{proof}

Together these results supply the representation, period-normalization
and scalar prerequisites of
Propositions~\ref{jacquet-original-budget-transport} and
\ref{jacquet-original-scale-transport}. They give the original local
leading-minor Jacquet estimate, including both character scales.
They do not prove that a prescribed global period map is nonzero:
the constructed local map is nonzero precisely when its original
period form is nonzero. The weighted orbital bound, rational-return
assembly and arithmetic integral theta maps remain separate.
